\chapter{Introduction}
\label{chp:intro}

The emergence of \ac{GenAI} and its most prominent representatives, \acp{LLM}, marked a technological inflection point. Over the past few years, this self-accelerating development has catalysed a profound transformation across a wide range of industries, changing the way knowledge is aggregated, processed and applied \cite{zhao2026surveylargelanguagemodels}. Even historically conservative and heavily regulated sectors such as accounting and tax consultancy are no exception. Law and consultancy firms face pressure on margins and fees, compelling them to adopt these technologies not as incremental aids but as strategic imperatives for operational efficiency \cite{mckinsey_genai_2023}. At the centre of this development lies the ambition to provide automated, intelligent support for time-consuming cognitive routine tasks, such as drafting tax memoranda or conducting initial legal assessments of complex cases \cite{lastres2015rebooting,waltl2018semantic,wind2026steuerllm}.

Yet the regulatory and technical demands of tax consultancy impose constraints that fundamentally distinguish it from the creative text-generation tasks at which \acp{LLM} excel. Legal and tax advice requires absolute factual accuracy and normative precision \cite{Larenz1995_Methodenlehre}. An erroneous statement can trigger extensive financial liability and regulatory consequences for both the adviser and the client \cite{Birk2023_Steuerrecht}. However, standard \acp{LLM} are vulnerable to hallucination\footnote{The confident generation of plausible but factually incorrect content.} and cannot reliably apply highly specialised or up-to-date expert knowledge, such as a firm's proprietary interpretations of current legislation \cite{Ji_2023_hallucination, zhao2026surveylargelanguagemodels}. Since continuous fine-tuning of the models to update knowledge is too inefficient and inflexible in industrial practice, alternative system architectures are required \cite{petroni_2019_language,barnett2024seven,lewis_rag_2020,zhao2026surveylargelanguagemodels}.

Against this background and in the absence of any market-available \ac{AI} solution that simultaneously satisfies the stringent requirements for legally valid advice and adaptively integrates proprietary in-house knowledge, \ac{PwC} developed the \acf{NEO}. \ac{NEO} is based on the paradigm of \acf{RAG}, a design that overcomes the limitations of a static parametric memory by externalising factual knowledge into an external corpus.  When a user makes a request, an \ac{IR} step first identifies the most relevant passages from the underlying corpus and injects them as grounding context into the \ac{LLM}'s generation process, thereby anchoring its output in source material \cite{lewis_rag_2020}.

While the discourse about \ac{LLM} hallucinations dominates the generative \ac{AI} landscape \cite{Ji_2023_hallucination}, this thesis shifts attention to the critical upstream bottleneck, the retrieval process itself. The quality of a \ac{RAG} system such as \ac{NEO} is bounded by the quality of its document search, if the retrieval algorithm fails to find the correct information, even the best generator will inevitably produce inadequate or misleading output \cite{barnett2024seven,reuter_towards_2025,pipitone2024legalbench}. Currently, \ac{NEO} uses a standard commercial embedding model (OpenAI's \texttt{text-embedding-3-large}) for dense retrieval. This universal embedder is widely adopted across the industry and trained to process a broad spectrum of texts, from poetry to recipes to news articles, a generality that is simultaneously its strength and its potential liability \cite{openai_embeddings_2024}.

Nevertheless, the German tax law domain poses fundamentally different requirements. Tax law texts are characterised by lexical homogeneity, formulaic phrasing and fine-grained semantic nuances, subtle distinctions in meaning that carry significant legal consequences \cite{van_opijnen_concept_2017,mattila_comparative_2013,busse_1992_recht}. It is a well-established finding in \ac{NLP} that universal embedding models suffer from a domain gap in such specialised settings \cite{gururangan_dont_2020,Lee_2019_BioBERT,chalkidis-etal-2020-legal}. Their representation spaces undergo a form of geometric collapse in which, although the model can separate tax texts from general-language documents, it fails to resolve the intra-domain distinctions that determine retrieval relevance \cite{ethayarajh-2019-contextual,wang-isola-2020-contrastive,bönisch2025finding}.

This problem is reinforced by the fact that in specialised domains such as German tax law, annotated reference datasets, the ground truth against which retrieval quality would conventionally be measured, are generally not available and their creation by domain experts is prohibitively expensive and time-consuming \cite{wrzalik_gerdalir_2021,pipitone2024legalbench,wind2026steuerllm}. In practice, this makes it considerably more difficult to objectively evaluate whether a universal embedding model is at all suitable for a given use case and how it might be systematically optimised or replaced \cite{ethayarajh_utility_2020,Sculley2018WinnersCO}. 

The present work addresses this gap. Rather than relying exclusively on costly, manually labelled evaluation datasets, this thesis approaches the problem from the perspective of vector space geometry. The aim is twofold. Firstly, it seeks to investigate the extent to which established, universal embedding models are suitable for \ac{RAG}-based retrieval in German tax law. Secondly, it examines whether the analysis of intrinsic geometric properties of embedding spaces can serve as a reliable predictor of actual, extrinsic retrieval performance. Appropriate quantitative metrics are applied to assess whether targeted interventions, ranging from domain-specific fine-tuning to resource-efficient post-processing transformations, can improve search quality measurably. From this objective, the central research question of this thesis emerges:

\begin{quote}
    How can the suitability of embedding models for the German tax law domain be systematically evaluated and optimised?
\end{quote}

To address this question in a structured manner, the investigation is decomposed into three guiding sub-questions:

\begin{enumerate}
    \item Geometric Characterisation and Response: What geometric properties do standard embeddings exhibit in the German tax law domain and how do these properties respond to different optimisation strategies?
    %\item Empirical Characterisation: What geometric and task-related properties do standard embeddings exhibit in the highly homogeneous German tax law domain?
    \item Predictive Validity: To what extent do intrinsic geometric properties correlate with actual retrieval quality and can they therefore serve as a reliable diagnostic signal in the absence of gold-standard evaluation data?
    %\item Targeted Optimisation: Through which specific interventions, such as domain-specific fine-tuning or post-hoc geometric transformations, can the embedding representations be optimised for this domain and do such corrections translate into measurably improved retrieval performance?
    \item Retrieval Improvement: Which of the tested optimisation strategies, such as domain-specific fine-tuning, model substitution, or post-hoc geometric transformations, yields the greatest measurable improvement in task-related retrieval quality for this domain?
\end{enumerate}

To delineate the scope of this study with precision, the work deliberately focuses exclusively on the dense retrieval step of the \ac{RAG} pipeline. While overall system performance could be enhanced through various complementary pipeline optimisations, such as hybrid search strategies, query expansion or neural re-ranking \cite{gao_2024_rag}, these extensions are intentionally excluded to avoid introducing confounding variables that would obscure the intrinsic capabilities of the embedding models under investigation. By holding the surrounding pipeline architecture constant, this thesis aims to derive fundamental, transferable recommendations for the selection,  efficient optimisation and performance-investigation of embedding models in specialised industrial domains.